\chapter{The IITG MA102 Examination Dataset and Practical Challenges}
\label{ch:dataset}

The foundation of our experimental work is the \textbf{MA102 Examination Dataset}, which we built using real undergraduate answer booklets from the Department of Mathematics at the Indian Institute of Technology Guwahati (IIT Guwahati). Rather than testing on clean synthetic equations rendered from arXiv or isolated CROHME crops, we wanted our models to face the messy realities of actual student exams: dense multi-line matrix reductions, cursive English proofs, slanted handwriting, uneven pen pressure, and heavy scratch-outs. This chapter explains how the dataset was collected, categorizes the types of mathematics students write, discusses the main handwriting degradations we encountered, and details how we split the data and built our ground-truth benchmarks.

\section{Course Background and Dataset Details}
\label{sec:dataset_demographics}

\textbf{MA102: Mathematics-II} is a mandatory first-year course for all undergraduate B.Tech students at IIT Guwahati. The syllabus covers two core areas:
\begin{enumerate}
    \item \textbf{Linear Algebra:} Vector spaces, subspaces, linear independence, basis and dimension, linear maps, matrix representations, rank-nullity theorem, systems of linear equations, eigenvalues, eigenvectors, and diagonalizability.
    \item \textbf{Multivariable Calculus:} Functions of several variables, partial derivatives, gradients, double and triple integrals, and fundamental vector calculus theorems.
\end{enumerate}

\subsection{Physical Booklets and Scanning Parameters}
During examinations, each candidate receives a standard unruled 14-page physical booklet. Once the exam is finished, the booklets are unbound and digitized using high-speed flatbed scanners at the Computer \& Communication Centre (CCC) at IIT Guwahati. Table \ref{tab:dataset_specs} lists the exact technical specifications of this digitized corpus.

\begin{table}[htbp]
\centering
\caption{IIT Guwahati MA102 Examination Dataset Specifications}
\label{tab:dataset_specs}
\vspace{0.2cm}
\small
\begin{tabular}{ll}
\toprule
\textbf{Parameter} & \textbf{Specification / Value} \\
\midrule
Academic Institution & Indian Institute of Technology Guwahati \\
Department & Department of Mathematics \\
Target Course & MA102 (Mathematics-II: Linear Algebra \& Calculus) \\
Student Cohort Size & 939 Enrolled Students \\
Booklet Length & 14 Physical Pages per Candidate \\
\textbf{Total Scanned Volume} & \textbf{13,146 Scanned Pages} (Single Exam Cycle) \\
Scanner Modality & Industrial Flatbed Optical Scanner \\
Optical Resolution & 200--300 Dots Per Inch (DPI) \\
Image Dimensions & $H = 2338 \times W = 1653$ pixels (Standard A4 Scan) \\
Color Space & 24-bit sRGB and 8-bit Grayscale Intensity \\
File Formats & JPEG / PNG \\
Paper Substrate & Unruled White Bond Paper \\
\bottomrule
\end{tabular}
\end{table}

\subsection{Printed Header versus Handwritten Ink}
Every scanned page is divided into two distinct zones:
\begin{itemize}
    \item \textbf{Printed Institutional Header ($y \le 0.20 H$):} The top $15\%$ to $20\%$ of the page contains standard printed department information: the course name (`MA102'), student metadata fields (Name and Roll Number), a question label (*"Space for answering Question-X only"*), and a solid horizontal printed dividing rule.
    \item \textbf{Handwritten Answer Area ($y > 0.20 H$):} The remaining $80\%$ of the page is completely blank and unruled. Students write their derivations here without any guidelines, leading to drifting lines, margin notes, and irregular line spacing.
\end{itemize}

\section{Types of Mathematical Problems in the Corpus}
\label{sec:math_taxonomy}

The student solutions in MA102 fall into three main structural types:

\subsection{2D Matrices and Gaussian Elimination}
In linear algebra problems, students frequently row-reduce large matrices to find ranks or invert systems. A clear example from student script \texttt{page-210103096-01} features a $4 \times 4$ real matrix containing negative signs and decimal fractions:
\begin{equation}
\label{eq:matrix_example}
\mathbf{A} = \begin{bmatrix}
-1.0 & 8.5 & -2.0 & 0.5 \\
4.5 & 2.5 & -3.0 & 3.0 \\
-0.5 & -1.5 & 3.0 & 1.0 \\
0.0 & 2.0 & -1.0 & 4.0
\end{bmatrix} \in M_{4 \times 4}(\mathbb{R})
\end{equation}

Transcribing these matrices accurately requires several steps:
\begin{enumerate}
    \item \textbf{Finding Delimiters:} Detecting tall outer brackets ($[ \dots ]$ or $| \dots |$) that span several vertical lines of handwriting.
    \item \textbf{Aligning the 2D Grid:} Grouping numbers into their correct rows and columns without accidentally merging adjacent numbers or dropping minus signs ($-$) and decimal points ($.$).
    \item \textbf{Parsing Row Operations:} Reading the handwritten transition arrows between steps, such as $R_2 \to R_2 + 4.5 R_1$ and $R_3 \to R_3 - \frac{1}{2} R_1$.
\end{enumerate}

\subsection{Deductive Vector Subspace Proofs}
When students prove whether a set $S \subset M_n(\mathbb{R})$ or $S \subset \mathbb{R}^n$ forms a subspace, they write step-by-step proofs testing three conditions:
\begin{enumerate}
    \item \textbf{Zero Vector Check:} $\mathbf{0} \in S \iff \begin{bmatrix} 0 & 0 \\ 0 & 0 \end{bmatrix} \in S$
    \item \textbf{Additive Closure:} $\forall \mathbf{u}, \mathbf{v} \in S \implies \mathbf{u} + \mathbf{v} \in S$
    \item \textbf{Scalar Multiplication Closure:} $\forall c \in \mathbb{R}, \ \forall \mathbf{u} \in S \implies c\mathbf{u} \in S$
\end{enumerate}
In these sections, students write in continuous cursive English (*"Let $\mathbf{u}, \mathbf{v} \in S$, then by definition of the subspace..."*). An OCR model must smoothly switch between reading English words and mathematical formulas without getting stuck in one mode.

\subsection{Linear Maps and Multivariable Calculus}
Students write definitions of linear maps $T: V \to W$, write out null spaces ($\ker(T) = \{ \mathbf{v} \in V \mid T(\mathbf{v}) = \mathbf{0} \}$), compute dimensions using the rank-nullity theorem ($\dim(V) = \text{rank}(T) + \text{nullity}(T)$), and solve multivariable calculus problems involving gradients ($\nabla f$), double integrals ($\iint_D f(x, y) \, dx dy$), and Greek variables ($\alpha, \beta, \lambda$). These problems mix nested subscripts with continuous prose.

\section{Real-World Handwriting and Physical Artifacts}
\label{sec:degradation_challenges}

Unlike clean benchmark datasets, real exam pages have many physical imperfections that complicate transcription:

\subsection{Handwriting Differences Across Students}
Across 939 students, penmanship varies wildly. Slant angles range from $-35^\circ$ to $+35^\circ$, and stroke thickness changes depending on whether the student used a thin ballpoint pen or a thick gel pen. When students write quickly under exam stress, pens sometimes run dry, turning thin strokes like minus signs or fraction bars into broken, disconnected dashes.

\subsection{Curved Lines and Baseline Drift}
Because the exam booklets are completely unruled, students cannot keep their handwriting perfectly horizontal. Baselines sag in the middle or tilt upward by $10^\circ$ to $15^\circ$. Simple horizontal projection cuts fail here, because the right half of an upper line often dips below the left half of the line beneath it.

\subsection{Interlocking Ascenders and Descenders}
A common challenge in handwritten math is vertical overlap between lines without the strokes actually touching. As shown in Figure \ref{fig:ascender_descender}, descenders like $y, g, p$ hang down right into the space where ascenders like $d, t, \lambda, \int$ reach up. A standard rectangular bounding box will slice right through these letters.

\begin{figure}[htbp]
\centering
\fbox{\parbox{0.85\textwidth}{\centering
\vspace{0.3cm}
\textbf{[Schematic Representation of Ascender-Descender Interlocking]}
\\[0.2cm]
$\text{Line } k: \quad \text{Let } y \text{ be a vector such that } p(y) = 0$ \quad (Descenders: $y, p$ extend down) \\
\vspace{0.1cm}
$\text{Line } k+1: \quad \text{Then } \lambda f(x) \text{ satisfies the theorem}$ \quad (Ascenders: $\lambda, f, t$ extend up) \\
\vspace{0.2cm}
\textit{Vertical bounding boxes overlap by $30$--$45\%$, requiring 1D Vertical IoU clustering.}
\vspace{0.3cm}
}}
\caption{Ascender-Descender vertical bounding overlap on unruled examination scripts.}
\label{fig:ascender_descender}
\end{figure}

\subsection{Crossed-Out Math and Scribbles}
When students spot an algebra mistake midway through an exam, they quickly scratch it out with heavy cross-hatching or diagonal strokes before starting over. Standard OCR models get confused by these overlapping lines and try to transcribe the scribbles as nonsense symbols. In our preprocessing, we detect regions with dense overlapping strokes so we can exclude them from the transcription.

\section{Dataset Partitioning and Ground-Truth Curation}
\label{sec:dataset_partitioning}

A common mistake in document AI experiments is \textbf{handwriting style leakage}. If crops from the same student are randomly scattered across training and test splits, the model simply memorizes that student's writing style, giving unrealistically high test scores.

\subsection{Student-Disjoint Partitioning Protocol}
To keep our benchmarks honest, we split the MA102 corpus strictly by \textbf{Student Roll Number}:
\begin{equation}
\mathcal{S}_{\text{total}} = \mathcal{S}_{\text{train}} \cup \mathcal{S}_{\text{val}} \cup \mathcal{S}_{\text{test}}, \quad \text{where } \mathcal{S}_i \cap \mathcal{S}_j = \emptyset \quad \forall i \neq j
\end{equation}
We divided the 939 students into three non-overlapping groups:
\begin{itemize}
    \item \textbf{Training Split ($\mathcal{S}_{\text{train}}$):} 750 students ($\approx 10,500$ scanned pages). Used to train our custom DenseNet-Coverage HMER models, tune projection parameters, and build token dictionaries.
    \item \textbf{Validation Split ($\mathcal{S}_{\text{val}}$):} 89 students ($\approx 1,246$ scanned pages). Used to tune line-grouping thresholds ($\theta_{\text{IoU}}$) and monitor training loss.
    \item \textbf{Test Split ($\mathcal{S}_{\text{test}}$):} 100 students ($\approx 1,400$ scanned pages). Kept completely untouched during all training and tuning.
\end{itemize}

\subsection{Ground-Truth Transcription Protocol and Benchmark Suite}
From the held-out test split $\mathcal{S}_{\text{test}}$, we selected a representative benchmark suite of authentic full pages capturing diverse mathematical challenges across undergraduate cohorts:
\begin{enumerate}
    \item \textbf{Student Roll 240101079 (Page 01):} Deductive vector subspace proof for symmetric matrices $W = \{ A \in M_{4 \times 4}(\mathbb{R}) \mid A^T = A \}$. Features mixed cursive English definitions and multi-step derivations checking all three vector space axioms. Used for our primary Multi-Band Slicing benchmark.
    \item \textbf{Student Roll 240122044 (Page 14):} Multi-step Gaussian elimination with row operations ($R_1 \to -R_1$, $R_2 \to R_2 - 3R_1$), where the lower canvas is unwritten, serving as our primary testbed for blank-paper hallucination loops.
    \item \textbf{Student Roll 210103096 (Pages 01 through 04):} Dense $4 \times 4$ decimal matrices, elementary row operations, coordinate equation systems, and multivariable calculus derivations.
    \item \textbf{Student Roll 240104105 (Pages 02 and 07):} Linear transformation null spaces ($\ker(T)$) and multivariable calculus double integrals.
\end{enumerate}

Table \ref{tab:benchmark_suite_summary} summarizes the characteristics of this benchmark suite.

\begin{table}[htbp]
\centering
\caption{Held-Out Benchmark Test Suite from MA102 Examination Corpus}
\label{tab:benchmark_suite_summary}
\vspace{0.2cm}
\small
\begin{tabular}{p{2.6cm}p{2.2cm}p{4.2cm}p{4.8cm}}
\toprule
\textbf{Student Roll No.} & \textbf{Exam Script} & \textbf{Mathematical Notation} & \textbf{Evaluation Role} \\
\midrule
Roll 240101079 & Page 01 & Subspace Proof: 3 Vector Axioms & Multi-Band Slicing Benchmark (88.24\% Accuracy) \\
\addlinespace
Roll 240122044 & Page 14 & Row Reduction \& Blank Canvas & Over-Helpful Tutor \& Blank Paper Loop Diagnosis \\
\addlinespace
Roll 210103096 & Pages 01--04 & $4 \times 4$ Decimal Matrix; Coordinates & Zero-Shot VLM Page-by-Page Baseline Analysis \\
\addlinespace
Roll 240104105 & Pages 02, 07 & Linear Maps ($\ker(T)$); Integrals & Modular Hybrid Dispatcher \& DAG Layout Analysis \\
\bottomrule
\end{tabular}
\end{table}

We hand-transcribed each benchmark page into \LaTeX\ under a strict \textbf{Verbatim Ink Protocol}:
\begin{enumerate}
    \item Every character, variable, operator, matrix entry, and punctuation mark is transcribed exactly as the student wrote it.
    \item \textbf{No Error Correction:} If a student wrote $-2.0$ instead of $+2.0$ in a row operation, our ground truth records $-2.0$. We never alter the student's work to match the professor's solution key.
    \item We compiled every ground-truth file with \texttt{pdflatex} and \texttt{tectonic} to ensure it compiles with zero syntax errors.
\end{enumerate}
This benchmark gives us an authentic ground truth for measuring EMR, TER, NED, and compilation rates in our experiments.
